\documentclass[11pt, english]{article}

\title{\textbf{COM3105 Advanced Algorithms} \\ Homework 4}
\author{Joffray Hargreaves}
\date{\vspace{-3em}} % Moves the date up or removes it completely

\usepackage[a4paper,
            bindingoffset=0in,
            left=1in,
            right=1in,
            top=1in,
            bottom=1in,
            footskip=.25in]{geometry}

\usepackage{amsthm} % For theorem environments
\usepackage{amsmath} % For advanced math typesetting
\usepackage{amssymb} % For \emptyset
\usepackage{bm} % For bold math symbols
\usepackage{setspace} % For line spacing
\setstretch{1.1} % Slightly increased line spacing for readability

\renewcommand{\labelitemi}{$\bullet$}
\renewcommand{\labelitemii}{$\circ$}
\renewcommand{\labelitemiii}{$\diamond$}

% --- Custom Environments ---

% Question environment: with bottom rule and marks aligned right
\newcounter{questionnum}
\newenvironment{question}[1][]{%
  \stepcounter{questionnum}%
  \par\noindent{\Large\bfseries Question \thequestionnum}%
  \ifx#1\empty\else\hfill{\normalsize\itshape(#1)}\fi%
  \par\vspace{0.25\baselineskip}\noindent\ignorespaces
}{%
  \par\vspace{1\baselineskip}%
}

% Answer environment: with top and bottom rule, no gap above heading
\newenvironment{answer}{%
  \par\noindent\hrule % top line
  \vspace{0.25\baselineskip}%
  \noindent{\Large\bfseries Answer}\par\vspace{0.25\baselineskip}\noindent\ignorespaces
}{%
  \par\vspace{0.5\baselineskip}%
  \noindent\hrule\vspace{0.5\baselineskip}%
  \hrule\hrule
}


% --- Document Start ---
\begin{document}

\setlength{\parindent}{0pt}  % remove indentation
% \setlength{\parskip}{0.9em}    % add space between paragraphs

\maketitle

\begin{question}[5 points]
Consider a stream \( a_1, a_2, \dots, a_m \), where each \( a_i \in [n] \).

Design and analyze a streaming algorithm that, given a query 

\( S \subseteq [n] \), 

outputs an estimate of what fraction of items in the stream belong to \( S \) within additive error \( \epsilon \) with probability at least \( \delta \). 

The space used by the algorithm should only depend on \( \epsilon \) and \( \delta \). 

Note that \( S \) is given only at query time (not in advance).

\end{question}

\begin{answer}

The intuition behind my answer is to use reservoir sampling with $k$ samples to estimate the fraction of streamed elements in $S$.

\textbf{Algorithm:}


\begin{itemize}
    \item \textbf{Initialisation:}  
    Choose a sample size $k$ based on the parameters:
    \[
        k \ge \frac{1}{4\,\epsilon^{2}\,\delta}
    \]

    Where $\epsilon$ is the chosen additive error and $\delta$ is the probability of failure.
    Initialise an empty set (the reservoir) $R$ that stores $k$ items.
    Initialise a counter $C = 0$ to track the number of items seen so far.

    \item \textbf{Processing the stream:}  
    For each incoming element $a_i$, where $i$ is the index of the element in the stream (starting from 1):
      \begin{itemize}
        \item If $C < k$, insert $a_i$ into $R$, and increment $C$ by 1.
        \item Otherwise:
          \begin{itemize}
          \item Include $a_i$ in the reservoir $R$ with probability $k/i$.
          \item If included, replace a uniformly random element from $R$.
          \item Increment $C$ by 1.
          \end{itemize}
      \end{itemize}
    This is standard reservoir sampling.

    \item \textbf{Query:}  
    When the set $S$ is given, compute the estimate:
    \[
        \widehat{p} = \frac{1}{k} \sum_{x \in R} \mathbf{1}\{x \in S\},
    \]
    and output $\widehat{p}$ as the estimate of the true fraction.

\end{itemize}

\bigskip

\textbf{Expectation of the Estimator:}

Let $p$ be the true fraction of stream elements that belong to $S$:
\[
    p = \frac{1}{m} \, |\{ i : a_i \in S \}|
\]

After reservoir sampling has finished, the reservoir contains
$k$ uniformly random elements from the stream.  
For each $j = 1, \ldots, k$, define the indicator variable
\[
    X_j = \mathbf{1}\{ R_j \in S \},
\],
where $R_j$ is the $j$-th sampled item.

Because reservoir sampling picks each stream element with equal probability,
each $R_j$ is distributed exactly as a uniformly random stream element.
So,
\[
    \mathbb{P}(R_j \in S) = p,
\]
and therefore
\[
    \mathbb{E}[X_j] = 1 \cdot p + 0 \cdot (1-p) = p
\]

The estimator for the fraction is
\[
    \widehat{p}
    = \frac{1}{k} \sum_{j=1}^k X_j
\]

Using linearity of expectation,
\begin{align*}
    \mathbb{E}[\widehat{p}]
    &= 
    \mathbb{E}\!\left[ \frac{1}{k} \sum_{j=1}^k X_j \right] \\[6pt]
    &= 
    \frac{1}{k} \sum_{j=1}^k \mathbb{E}[X_j] \\[6pt]
    &= 
    \frac{1}{k} \sum_{j=1}^k p \\[6pt]
    &= 
    \frac{1}{k} \cdot k p \\[6pt]
    &= p
\end{align*}

So the estimator has expectation equal to the true fraction:
\[
    \mathbb{E}[\widehat{p}] = p
\]

\bigskip

\textbf{Variance of the Estimator:}

Remember that each random variable $X_j$ is an indicator:
\[
    X_j =
    \begin{cases}
        1 & \text{if the $j$-th sampled element belongs to } S, \\
        0 & \text{otherwise}
    \end{cases}
\]
So,
\[
    \mathbb{P}(X_j = 1) = p
    \qquad\text{and}\qquad
    \mathbb{P}(X_j = 0) = 1 - p
\]

To get the variance of $X_j$ we can use the formula:
\[
    \mathrm{Var}(X_j)
    = \mathbb{E}[X_j^2] - (\mathbb{E}[X_j])^2
\]

Since $X_j^2 = X_j$ (as the indicators are either 0 or 1):
\[
    \mathbb{E}[X_j^2]
    = \mathbb{E}[X_j]
\]

And from the expectation calculation above we have:
\[
    \mathbb{E}[X_j] = p
\]

Putting this together:
\[
    \mathrm{Var}(X_j)
    = p - p^2
    = p(1-p).
\]

Because $p(1-p)$ is maximized when $p=\tfrac12$, we know that:
\[
    \mathrm{Var}(X_j) = p(1-p) \le \frac14
\]

\bigskip

Remembering that the estimator is the mean of the $X_j$ in the reservoir:
\[
    \widehat{p} = \frac{1}{k} \sum_{j=1}^k X_j
\]

Remembering the two properties of variance:
\begin{itemize}
    \item For any random variable $Y$ and constant $c$,
    \[
        \mathrm{Var}(cY) = c^2 \, \mathrm{Var}(Y)
    \]
    \item For independent random variables $Y_1, Y_2$,
    \[
        \mathrm{Var}(Y_1 + Y_2 )
        =
        \mathrm{Var}(Y_1) + \mathrm{Var}(Y_2)
    \]
\end{itemize}


As the $X_j$ are independent, we can use these properties to compute the variance of $\widehat{p}$:
\[
    \mathrm{Var}(\widehat{p})
    =
    \mathrm{Var}\!\left( \frac{1}{k} \sum_{j=1}^k X_j \right)
    =
    \frac{1}{k^2} \sum_{j=1}^k \mathrm{Var}(X_j)
\]

Using $\mathrm{Var}(X_j) \le 1/4$ for every $j$:
\[
    \mathrm{Var}(\widehat{p})
    \le
    \frac{1}{k^2} \cdot k \cdot \frac14
    =
    \frac{1}{4k}
\]
\newpage

\textbf{Quality Analysis Using Chebyshev's Inequality:}

Chebyshev's inquality is defined as:

\[
\mathbb{P}(|X - \mu| > \lambda) \leq \frac{\text{Var}(X)}{\lambda^2}
\]

Using Chebyshev’s inequality for our estimator $\widehat{p}$, we can substitute in our bounded variance (from above) and set $\lambda = \epsilon$:
\[
    \mathbb{P}\left( |\widehat{p} - p| > \epsilon \right)
    \le
    \frac{1}{4k\epsilon^2}
\]

To bound the probability of error, we set the right-hand side to be at most $\delta^\prime$:
\[
    \frac{1}{4k\epsilon^2} \le \delta^\prime
\]

Now we solve for $k$:
\[
    \frac{1}{4\epsilon^2 \delta^\prime} \le k
\]

Substituting this into the right-hand side:
\[
    \mathbb{P}\big(|\widehat{p} - p| > \epsilon \big)
    \le \frac{1}{4 \cdot \left(\frac{1}{4 \epsilon^2 \delta^\prime}\right) \cdot \epsilon^2 }
    = \frac{1}{ \frac{1}{\delta^\prime} } 
    = \delta^\prime
\]

Therefore, with this choice of \(k\) we have:
\[
    \mathbb{P}\big(|\widehat{p} - p| > \epsilon \big) \le \delta^\prime
\]

Which we can rewrite as:
\[
    \mathbb{P}\big(|\widehat{p} - p| \le \epsilon \big) \ge 1 - \delta^\prime
\]

Now we can set \(1 - \delta^\prime = \delta\) to bound the probability of error by \(\delta\):
\[
    \mathbb{P}\big(|\widehat{p} - p| \le \epsilon \big) \ge \delta
\]

This proves that the estimator $\widehat{p}$ is within additive error $\epsilon$ of the true fraction $p$ with probability at least $\delta$.


\bigskip

\textbf{Space Complexity Analysis:}

The algorithm stores:
\begin{itemize}
    \item A reservoir of size $k = \mathcal{O}\left(\frac{1}{4\epsilon^2 (1-\delta)}\right)$.
    \item Each stored item in $k$ is from $[n]$ and requires $\mathcal{O}(\log n)$ bits.
\end{itemize}

Therefore, the total space used by the algorithm is:
\[
    \mathcal{O}\!\left( k \log n \right)
    =
    \mathcal{O}\!\left( \frac{1}{\epsilon^2 (1-\delta )} \log n \right)
\]

This space depends only on $\epsilon$ and $\delta$ (and $\log n$, but this is not tuneable), as required.

\end{answer}

\vspace{1em}

\newpage

\begin{question}[5 points]
question here
\end{question}

\begin{answer}
answer here
\end{answer}

\end{document}



% Proof that reservoir sampling gives each element with equal probability, assuming that the stream length is larger than $k$:
% \begin{itemize}
%     \item Let $a_{m-1}$ be the $m-1$-th element in the stream, that has been processed.
%     \item $a_{m}$ is the next element to be processed.
%     \item From the algorithm definition above, we get:
%       \begin{itemize}
%         \item probability of $a_{m}$ being included in the reservoir = $k/m$
%         \item probability of not being included = $1 - k/m$
%       \end{itemize}
%     \item hence, by definition, the probability of $a_{m}$ being in the reservoir after processing $m$ elements is $k/m$.
%     \item For any other element $a_i$ where $i < m$, let $R_{m-1}$ be the reservoir after processing $m-1$ elements, and $R_{m}$ be the reservoir after processing $m$ elements:
%       \begin{itemize}
%         \item $\mathbb{P}[a_i \in R_{m}] = \mathbb{P}[a_i \in R_{m-1}] \cdot \mathbb{P}[a_{m} \text{ did not replace } a_i]$
%         \item simplify the left term:
%           \[
%               \mathbb{P}[a_i \in R_{m-1}] = \frac{k}{m-1}
%           \]
%         \item simplify the right term:
%         \begin{itemize}
%           \item $\mathbb{P}[a_{m} \text{ did not replace } a_i] = \mathbb{P}[a_{m} \text{did not replace any element in } R_{m-1}] + \mathbb{P}[a_{m} \text{ replaced an element in } R_{m-1} \text{ but not } a_i]$        
%         \end{itemize}
%       \end{itemize}

% \end{itemize}
    